% Prose for this counterexample.  Every region below is delimited by an
% environment that tools/build.py extracts; machine facts (ids, uid, dates,
% urls, enums) live in case.json.  See counterexamples/AGENTS.md.

\cxtitle{Quantum coupon collection: an alternating sum of inverses}

\begin{cxcredits}
\posedby{S.~Sra \citeyearpar{Sra2017QuantumCoupon}}
\foundby{GPT Pro}{2026-02-16}
\formalizedby{GPT-5.6 Pro, which re-derived the exact witnesses in 2026-08}
\auditedby{GPT-5.6 Pro}
\contributedby{S.~Sra}
\end{cxcredits}

\begin{cxcontext}
In the classical coupon collector's problem with sampling probabilities
$x_1,\ldots,x_n>0$, the expected waiting time is the alternating sum
\[
T_n(x_1,\ldots,x_n)=\sum_{k=1}^n(-1)^{k+1}
\sum_{1\le i_1<\cdots<i_k\le n}\frac{1}{x_{i_1}+\cdots+x_{i_k}},
\]
whose positivity can be seen through a nice trick using the fact that $t\mapsto1/t$ is completely monotone on $(0,\infty)$, so the Laplace-transform representation makes the alternating sum an integral of a nonnegative quantity \cite{FlajoletGardyThimonier1992}.  The statement quoted below asks whether the same holds with the positive numbers replaced by positive definite matrices. % and the reciprocal by the matrix inverse.

Let $\mathbf{S}_{++}^{d}$ be the cone of real symmetric positive definite $d\times d$ matrices.
% For a nonempty $S\subseteq[n]$ we abbreviate $X_S:=\sum_{i\in S}X_i$, so that the source's double sum is
% \[
% Q_n(X_1,\ldots,X_n)=\sum_{\emptyset\ne S\subseteq[n]}(-1)^{|S|-1}X_S^{-1}.
% \]
% We also write $\Delta_n:=\tr Q_n$ for its scalar trace transform; since a positive definite matrix has positive trace, the conjecture implies $\Delta_n>0$.

The cone-valued form of the scalar mechanism is available: Theorem~7 of Niculescu and Sra \cite{NiculescuSra2023} states that every completely monotone $f$ on a cone $\mathcal C$ satisfies the inclusion--exclusion inequality $\sum_i f(x_i)-\sum_{i<j}f(x_i+x_j)+\cdots\ge0$ for all $x_1,\ldots,x_n\in\mathcal C$.  Applying it to $\mathcal C=\mathbf{S}_{++}^{d}$ and $f_v(X)=v^{\mathsf T}X^{-1}v$ would settle the conjecture for every $v$, and applying it to $X\mapsto\tr(X^{-1})$ would settle the trace form.  The refutation below therefore also shows that neither function is completely monotone on $\mathbf{S}_{++}^{d}$ for $d\ge3$ --- for $d=1$ both reduce to $1/t$, which is.  A matrix function that does satisfy the whole alternating string on $\mathbf{S}_{++}^{d}$ is $X\mapsto(\det X)^{-\rho}$, for $\rho\in\{0,\tfrac12,1,\tfrac32,\ldots\}\cup[(d-1)/2,\infty)$ \cite{NiculescuSra2023}.
\end{cxcontext}

% ---------------- result: quantum-coupon-collector ----------------

\begin{cxsource}{quantum-coupon-collector}
S.~Sra, MathOverflow question 263833 \citeyearpar{Sra2017QuantumCoupon}
\end{cxsource}

\begin{cxstatement}{quantum-coupon-collector}
Here we take symmetric positive definite matrices $X_1,\ldots,X_n\in\mathbf{S}_{++}^d$, and consider the sum
\[
Q_n(X_1,\ldots,X_n):=\sum_{k=1}^n(-1)^{k+1}
\sum_{1\le i_1<\cdots<i_k\le n}\left(X_{i_1}+\cdots+X_{i_k}\right)^{-1}.
\]

\textbf{Conjecture.} For $n\ge1$, and $X_1,\ldots,X_n\in\mathbf{S}_{++}^d$, we have $Q_n(X_1,\ldots,X_n)\succ0$.

\emph{Note.} The cases $n=1,2,3$ are trivially true.  I've tried $n=4$ numerically, and it seems to hold.  (Also, by CCP, positivity immediately holds for simultaneously diagonalizable matrices.)
\end{cxstatement}

\begin{cxsummary}{quantum-coupon-collector}
The matrix coupon-collector sum $\sum_{\emptyset\ne S}(-1)^{|S|-1}X_S^{-1}$ was conjectured positive definite for every order $n$ and every dimension $d$.
\end{cxsummary}

\begin{cxcertificate}{quantum-coupon-collector}
Exact rational $3\times3$ matrices with $n=6$ and the integer vector $(4,1,3)^{\mathsf T}$ give a quadratic form strictly between $-96$ and $-95$.
\end{cxcertificate}

\begin{cxrefutation}
The conjecture holds for a tuple of simultaneously diagonalizable matrices, where it is the classical identity coordinate by coordinate, and the source records that $n\le3$ is trivial and that $n=4$ was checked numerically.  Lemma~\ref{lem:qcc-facet} and Theorem~\ref{thm:qcc-positive-five} below upgrade that record: the conjecture is a theorem for every $d$ and every $n\le5$.  It fails at $n=6$.

\begin{lemma}[Facet averaging]\label{lem:qcc-facet}
Let $A_1,\ldots,A_m\in\mathbf{S}_{++}^{d}$, where $m\ge2$, and put $A=\sum_{j=1}^{m}A_j$.  Then
\[
A^{-1}\preceq\frac{m-1}{m^2}\sum_{j=1}^{m}(A-A_j)^{-1}.
\]
\end{lemma}
\begin{proof}
Set $B_j=A-A_j$, which is positive definite because it is a sum of the remaining $A_i$.  Since $\frac1m\sum_{j=1}^{m}B_j=\frac{m-1}{m}A$ and the inversion map is operator convex on $\mathbf{S}_{++}^{d}$,
\[
\frac{m}{m-1}A^{-1}
=\Bigl(\frac1m\sum_{j=1}^{m}B_j\Bigr)^{-1}
\preceq\frac1m\sum_{j=1}^{m}B_j^{-1}.
\]
Rearranging gives the claim.
\end{proof}

For a nonempty $S\subseteq[n]$ we abbreviate $X_S:=\sum_{i\in S}X_i$; then, for fixed $n$, introduce the layer sums
\[
E_k:=\sum_{\substack{S\subseteq[n]\\|S|=k}}X_S^{-1},\qquad1\le k\le n,
\]
each of which is positive definite.  Applying Lemma~\ref{lem:qcc-facet} inside each $(k+1)$-subset $T$, with the $A_j$ taken to be the $X_i$ for $i\in T$, gives $X_T^{-1}\preceq\frac{k}{(k+1)^2}\sum_{i\in T}X_{T\setminus\{i\}}^{-1}$.  Summing over all $T$ of size $k+1$, and noting that each fixed $k$-subset arises as a facet of exactly $n-k$ of them, yields
\begin{equation}\label{eq:qcc-layer}
E_{k+1}\preceq\frac{k(n-k)}{(k+1)^2}E_k,\qquad1\le k<n.
\end{equation}

\begin{theorem}[\statusproved: positivity through five variables]\label{thm:qcc-positive-five}
For every $d\ge1$, every $1\le n\le5$, and every $X_1,\ldots,X_n\in\mathbf{S}_{++}^{d}$, one has $Q_n(X_1,\ldots,X_n)\succ0$.
\end{theorem}
\begin{proof}
$n=1$ is immediate.  Write $Q_n=E_1-E_2+E_3-\cdots+(-1)^{n-1}E_n$ and apply \eqref{eq:qcc-layer} to obtain:
\[
\begin{array}{rcl}
n=2:&E_2\preceq\frac14E_1,&Q_2\succeq\frac34E_1,\\[0.35em]
n=3:&E_2\preceq\frac12E_1,&Q_3\succeq\frac12E_1+E_3,\\[0.35em]
n=4:&E_2\preceq\frac34E_1,\quad E_4\preceq\frac3{16}E_3,
&Q_4\succeq\frac14E_1+\frac{13}{16}E_3,\\[0.35em]
n=5:&E_2\preceq E_1,\quad E_4\preceq\frac38E_3,
&Q_5\succeq\frac58E_3+E_5.
\end{array}
\]
Each displayed lower bound is clearly positive definite.
\end{proof}

\begin{theorem}[\statusfalse: quantum coupon collection]\label{thm:quantum-coupon-collector}
The conjecture is false: there are $d=3$, $n=6$ and $X_1,\ldots,X_6\in\mathbf{S}_{++}^{3}$ with $Q_6(X_1,\ldots,X_6)\not\succeq0$, so in particular $Q_6\succ0$ fails.
\end{theorem}
\begin{proof}
Take $d=3$, $n=6$, $\varepsilon=1/100$, and
\[
X_i=w_iu_iu_i^{\mathsf T}+\varepsilon I_3,
\]
where
\[
\begin{array}{c|c|c}
i&u_i^{\mathsf T}&w_i\\\hline
1&(-1,2,1)&10\\
2&(0,-3,1)&10\\
3&(-3,-2,0)&100\\
4&(2,-2,-2)&100\\
5&(-1,0,2)&100\\
6&(-1,3,0)&100
\end{array}
\]
Each $X_i$ is strictly positive definite because $X_i\succeq\varepsilon I_3\succ0$, so the hypotheses are met exactly as posed.  The example is genuinely noncommutative --- the $(1,2)$ entry of $[X_1,X_2]$ equals $-1500$ --- and therefore lies outside the simultaneously diagonalizable regime in which the classical identity applies.  For $v=(4,1,3)^{\mathsf T}$, exact rational evaluation of the sixty-three nonempty-subset terms gives
\[
-96<v^{\mathsf T}Q_6(X_1,\ldots,X_6)v<-95.
\]
The accompanying verifier performs this computation using only rational arithmetic. Thus, $Q_6$ has a negative quadratic form and is not positive semidefinite, let alone positive definite.
\end{proof}

The failure is not confined to the matrix ordering: the weaker scalar consequence $\tr Q_n>0$ fails as well; see  Appendix~\ref{sec:qcctrace} for more details and commentary.

% \begin{remark}[Scope]
% What is refuted is the conjecture as posed, universally quantified over $n$ and $d$.  The exact record is now: for every $d$, the conjecture is true for $n\le5$ (Theorem~\ref{thm:qcc-positive-five}) and false at $n=6$ (Theorem~\ref{thm:quantum-coupon-collector}).  For the trace form the record is coarser --- true for $n\le5$, false at $n=10$, undecided in between; see Appendix~\ref{sec:qcctrace}.  Nothing here contradicts the simultaneously diagonalizable case, which remains an identity.
% \end{remark}
\end{cxrefutation}


%%% Local Variables:
%%% mode: LaTeX
%%% TeX-master: "../../tex/main"
%%% End:
